\documentclass{article}
\usepackage{amsmath, amssymb}
\usepackage{graphicx}

\begin{document}

\title{Understanding the Rendering Equation as a Convolution}
\author{}
\date{}
\maketitle

\section{Introduction}
The rendering equation is the foundation of physically based rendering (PBR). It models how light interacts with surfaces and is mathematically expressed as:

\begin{equation}
L_o(\omega_o) = \int_{\Omega} f_r(\omega_o, \omega_i) L_i(\omega_i) (\omega_i \cdot n) d\omega_i.
\end{equation}

This equation states that the outgoing radiance \( L_o(\omega_o) \) at a point is the integral of the incoming radiance \( L_i(\omega_i) \) weighted by the surface's Bidirectional Reflectance Distribution Function (BRDF) \( f_r(\omega_o, \omega_i) \) and a cosine term \( (\omega_i \cdot n) \), which accounts for the angle of incidence.

\section{Interpreting the Rendering Equation as a Convolution}
Mathematically, convolution is defined as:

\begin{equation}
(f * g)(t) = \int f(\tau) g(t - \tau) d\tau.
\end{equation}

Comparing this with the rendering equation, we can interpret it as a convolution where:

\begin{itemize}
    \item \( L_i(\omega_i) \) represents the incoming function (irradiance),
    \item \( f_r(\omega_o, \omega_i) \cdot (\omega_i \cdot n) \) acts as the convolution kernel (filter),
    \item The result is the outgoing radiance \( L_o(\omega_o) \), a transformed version of \( L_i \).
\end{itemize}

Thus, the rendering equation describes how **incoming light is filtered by the surface BRDF**, much like a signal being processed by a convolution filter.

\section{Importance Sampling and PDFs in Monte Carlo Integration}
Since direct computation of the integral is intractable, Monte Carlo integration is used:

\begin{equation}
L_o(\omega_o) \approx \frac{1}{N} \sum_{i=1}^{N} \frac{f_r(\omega_o, \omega_i) L_i(\omega_i) (\omega_i \cdot n)}{p(\omega_i)}.
\end{equation}

The Probability Density Function (PDF), \( p(\omega_i) \), determines which directions are sampled more frequently. By choosing \( p(\omega_i) \) proportional to \( \cos\theta \), we optimize the sampling efficiency, reducing noise while maintaining accuracy.

\section{Conclusion}
This interpretation lays the foundation for understanding **BxDFs** as convolution filters that shape light reflections, and **Image-Based Lighting (IBL)** as an extension of the rendering equation. The next sections will explore these concepts in detail.

\end{document}
